\scanpage{37}
\mspage{14}

\deleted{\unsure{NB}}\quad \uline{Remarks}\quad
1) With hypotheses of ``non-wildness'' on the initial datum
\[
 (X,(X_i)_{i\in I},(B_\alpha,\mathcal E_\alpha)_{\alpha\in\Lambda}),
\]
it should be true that there exists a conical structure. If the stratification
\deleted{\unsure{X}} is equisingular along every
\[
 X_i^\circ\mathrel{\mathop=^{\mathrm{def}}}
 X_i\setminus\bigcup_{\substack{j<i\\j\le i}}X_j,
\]
then there should \(=\) be an equisingular conical structure
(with the obvious ad-hoc definition) --- it suffices to know how to construct
an equisingular tube around the \(X_{I_0}\) \(\cdots\)

2) A priori, if one has a \inserted{strictly} increasing sequence
of subsets
\[
 \phi\subset I_0\subset I_1-\, -\subset I_{n-1}\subset I_n=I,
\]
with \deleted{\(I_i\) consisting of minimal elements}
\inserted{\(\forall i\)} \(I_i\setminus I_{i-1}\) consisting of minimal elements
of \(I\setminus I_{i-1}\), then one could define
\inserted{by induction} a notion of conical structure
for this ``filtration'' of \(I\). But one
can show that this notion does not depend
in fact, essentially, on the filtration
chosen.

3) Instead of a family of \((B_\alpha,\mathcal E_\alpha)_{\alpha\in\Lambda}\),
one can suppose given on \(X\) a
cylindrical stratification \deleted{of type} (the case envisaged
is that where the ordered set which \deleted{\unsure{poin|carries}}
\unsure{describes} the cylindrical stratification is
\(\mathcal P(\Lambda)^\circ\).

\margin{36}
